\documentclass[a4paper,10pt]{ltjsarticle}

\usepackage{luatexja-fontspec}
\usepackage{amsmath, amssymb}
\usepackage{geometry}
\geometry{margin=20mm}

\usepackage{hyperref}
\usepackage{setspace}
\setstretch{0.95}

\title{作るか、選ぶか?\\教員による提示方法が教育動画の価値認識に与える影響\\
\large Create or Curate?\\ How Instructor Framing Shapes the Value of Educational Video\\
1022179 西野理希 Riki Nishino\\
指導教員：Ian Frank}
\author{}
\date{}

\begin{document}

\maketitle

\section{研究背景}

近年、大学教育においてDX（デジタルトランスフォーメーション）の推進やオンライン学習環境の普及に伴い、教育動画の活用が広がっている。Noetelら（2021）は、高等教育における教育動画に関するシステマティックレビューおよびメタ分析を行い、教育動画が学習成果の向上に寄与することを示している。このことから、教育動画は大学教育における有効な学習資源として広く活用されている。
一方で、教育動画を新たに制作するためには、時間や技術に加え、一定のデジタル能力が求められる。Liesa-Orúsら（2023）は、大学教員のデジタル能力には大きな差が存在することを報告しており、すべての教員が独自の動画教材を制作できるわけではないことを指摘している。そのため、教育現場では、動画を新たに制作するだけでなく、既存の動画教材を適切に選定し活用することも重要な実践となる。
このような状況の中で、教育動画そのものだけでなく、教員が動画をどのように位置付けて提示するかにも注目する必要がある。同一の動画であっても、「教員が制作した教材」と説明される場合と、「教員が厳選した教材」と説明される場合とでは、学習者の受け止め方が異なる可能性がある。しかし、このような教員による提示方法が教育動画の評価に与える影響については、十分に検討されていない。

\section{研究目的}

本研究では、教育動画に付与される教員由来の情報が、学習者による動画の評価にどのような影響を与えるかを明らかにすることを目的とする。
具体的には、同一の動画教材を用いながら、「教員が制作した教材」と説明する場合と、「教員が厳選・キュレーションした教材」と説明する場合とで、学習者が感じる教育的価値、有用性、信頼性などの評価に違いが生じるかを検討する。
これにより、教育におけるキュレーションという概念がどの程度有効であるかを明らかにすることを目指す。

\section{研究課題（Research Question）}

教育における「キュレーション」という概念はどの程度有効なのか。より具体的には、同一の教育動画に対して異なる提示方法を用いた場合、学習者による教育動画の評価にどのような違いが生じるのかを明らかにする。

\section{関連研究}

教育動画の活用に関する研究では、教員は動画を新たに制作するだけでなく、既存の動画を選定して授業に取り入れることも一般的な教育実践であることが報告されている。Shoufan and Mohamed（2022）は、教育におけるYouTube活用に関する研究を整理し、教員は目的に応じて動画を制作あるいは選定して活用していることを示している。また、Deschaine and Sharma（2015）は、デジタルキュレーションを教育目的に沿って学習資源を選定・整理・共有する活動として位置付けており、Flintoffら（2014）も、教育・研究・専門能力開発を支える実践としてデジタルキュレーションの重要性を論じている。
一方、コミュニケーション研究では、同一の情報であっても、情報源に対する認識が受け手の評価に影響することが知られている。Hovland and Weiss（1951）は、情報源の信頼性が情報の説得力や受容に影響することを示した。
さらに、近年のオンライン教育研究では、教材内容が同一であっても、教材の提示方法や教員に関する情報によって学習者の評価が変化することが報告されている。Ramlatchan and Watson（2020）は、オンライン講義動画の提示方法のみを変更した実験により、提示方法の違いが学習者による教員の信頼性や心理的近接性の評価に影響することを示した。また、Bateman and Schmidt-Borcherding（2018）は、教育動画では映像や音声、レイアウトなど様々なシグナルが学習者の解釈や評価に影響すると論じている。
しかし、同一の教育動画を用いながら、「教員が制作した教材」と「教員が厳選した教材」という説明のみを操作し、その違いが学習者による教育動画の評価に与える影響を検討した研究はほとんど見られない。


\section{研究方法}

本研究では、同一の教育動画を用いた小規模な実験を実施する。
参加者は以下の3条件のいずれかに割り当てられる。

\begin{itemize}
  \item 制作条件：「この動画は担当教員がこの授業のために制作した動画です」と提示する。
  \item キュレーション条件：「この動画は担当教員が授業のために厳選した動画です」と提示する。
  \item 中立条件：動画の由来に関する追加説明を行わない。
\end{itemize}

すべての条件で同一の動画教材を使用する。
動画視聴後には理解度テストおよびアンケート調査を実施する。アンケートでは、教育的価値、有用性、信頼性、満足度、学習意欲などを測定する。
得られた結果を条件間で比較し、教員による提示方法が学習者の教育動画に対する認識や評価にどのような影響を与えるかを分析する。

\section{期待される成果}

本研究は、教育動画そのものの品質や学習効果を比較するのではなく、教員による提示方法が学習者の認識に与える影響に着目する。
これにより、教育動画を新たに制作することが難しい環境においても、既存の動画を適切に選定し、教育的な文脈の中で提示することの価値を検討することができる。
また、本研究は、教育動画の活用において「作ること」だけでなく「選び、位置付けること」の教育的意義を示す知見を提供することを目指す。


\begin{thebibliography}{99}

\bibitem{noetel2021}
Noetel, M., et al. (2021).
Video improves learning in higher education: A systematic review.
\textit{Review of Educational Research}, 91(2), 204--236.

\bibitem{liesa2023}
Liesa-Orús, M., et al. (2023).
Digital competence of university teachers: A systematic review.
\textit{Education Sciences}, 13(5), 508.

\bibitem{shoufan2022}
Shoufan, A., \& Mohamed, F. (2022).
YouTube and Education: A Scoping Review.
\textit{IEEE Access}, 10, 125547--125592.

\bibitem{deschaine2015}
Deschaine, M. E., \& Sharma, R. (2015).
The five stages of digital curation: Supporting teaching and learning in online learning environments.
\textit{Journal of Online Learning and Teaching}, 11(3), 398--409.

\bibitem{flintoff2014}
Flintoff, K., Mellow, P., \& Clark, M. (2014).
Digital curation: Opportunities for learning, teaching, research and professional development.
In \textit{Proceedings of ASCILITE 2014}.

\bibitem{hovland1951}
Hovland, C. I., \& Weiss, W. (1951).
The influence of source credibility on communication effectiveness.
\textit{Public Opinion Quarterly}, 15(4), 635--650.

\bibitem{ramlatchan2020}
Ramlatchan, M., \& Watson, G. S. (2020).
Enhancing instructor credibility and immediacy in online multimedia designs.
\textit{Educational Technology Research and Development}, 68, 3239--3264.

\bibitem{bateman2018}
Bateman, J. A., \& Schmidt-Borcherding, F. (2018).
The communicative effectiveness of education videos:
Towards an empirically motivated multimodal account.
\textit{Multimodal Technologies and Interaction}, 2(3), 59.

\end{thebibliography}

\end{document}